%------------------------------------------------------------------------------------
%  SECTION 5.1 : FLOTTANTS CLASSIQUES VS TYPE D'ILLUSTRATION DÉDIÉ
%------------------------------------------------------------------------------------
\section{Flottants Standards et Environnement Dédié aux Illustrations}
\label{sec:figures_illustrations}

La restitution graphique d'expériences ou d'analyses conceptuelles s'articule autour de deux types de supports visuels :
les figures techniques (graphiques vectoriels, courbes, schémas de principe) d'une part, et les illustrations iconographiques (clichés photographiques, reproductions anatomiques, œuvres artistiques) d'autre part.
Mélanger ces deux catégories au sein d'une unique « table des figures » nuit à la clarté du document.

\subsection{Figure standard avec sous-figures coordonnées (\textsf{subcaption})}
\label{subsec:sous_figures}

Pour juxtaposer deux résultats comparatifs, le paquet moderne \textsf{subcaption} remplace avantageusement le paquet obsolète \texttt{subfig}.
Chaque sous-figure est isolée dans un environnement \texttt{subfigure} avec sa propre sous-légende indicée \texttt{(a)}, \texttt{(b)} :

\begin{figure}[htbp]
  \centering
  \begin{subfigure}[b]{0.46\textwidth}
    \centering
    \begin{tcolorbox}[colframe=secondaryColor, colback=white, boxrule=1pt, arc=4pt, width=\linewidth, halign=center]
      \vspace{0.8cm}
      {\large\sffamily\bfseries\color{secondaryColor} [ Schéma Technique A ]}\\[2mm]
      {\footnotesize Flux d'injection sous pression \qty{150}{\kilo\pascal}}
      \vspace{0.8cm}
    \end{tcolorbox}
    \caption{Régime d'injection stationnaire}
    \label{subfig:injection_stationnaire}
  \end{subfigure}
  \hfill
  \begin{subfigure}[b]{0.46\textwidth}
    \centering
    \begin{tcolorbox}[colframe=secondaryColor, colback=white, boxrule=1pt, arc=4pt, width=\linewidth, halign=center]
      \vspace{0.8cm}
      {\large\sffamily\bfseries\color{secondaryColor} [ Schéma Technique B ]}\\[2mm]
      {\footnotesize Propagation d'ondes à \qty{4.2}{\meter\per\second}}
      \vspace{0.8cm}
    \end{tcolorbox}
    \caption{Régime transitoire perturbé}
    \label{subfig:injection_transitoire}
  \end{subfigure}
  \caption{Comparaison des dynamiques d'écoulement sous deux régimes de fonctionnement}
  \label{fig:comparaison_regimes}
\end{figure}

Comme le démontre la~\cref{fig:comparaison_regimes}, et plus spécifiquement la~\cref{subfig:injection_stationnaire}, chaque sous-figure peut être référencée indépendamment dans le corps du texte grâce à \textsf{cleveref}.

\subsection{L'innovation de \texttt{ClasseGuillaume} : Le flottant \texttt{illustration}}
\label{subsec:flottant_illustration}

Pour répondre aux standards des grandes monographies et thèses universitaires, la classe déclare un nouveau type de flottant officiel via le moteur \textsf{tocbasic} de KOMA-Script :

\begin{verbatim}
\DeclareNewTOC[%
  type = illustration,
  float,
  name = Illustration,
  listname = {Liste des illustrations},
]{loi}
\end{verbatim}

Cet environnement s'utilise exactement comme une figure ordinaire, mais bénéficie d'une numérotation propre et alimente une \textbf{table des illustrations indépendante} (\texttt{\textbackslash listofillustrations}) située en fin d'ouvrage :

\begin{illustration}[htbp]
  \centering
  \begin{cadreFigure}[width=0.82\textwidth]
    \vspace{0.6cm}
    {\Large\sffamily\bfseries\color{primaryColor} Cliché Iconographique de Microscopie Électronique}\\[3mm]
    {\normalsize Grossissement $\times \num{25000}$ -- Structure cristalline des nanotubes de carbone}
    \vspace{0.6cm}
  \end{cadreFigure}
  \caption{Observation nanoscopique d'échantillons en champ sombre}
  \label{ill:microscopie_nanotubes}
\end{illustration}

L'\cref{ill:microscopie_nanotubes} est reconnue de manière transparente par \textsf{cleveref}, qui lui attribue automatiquement le préfixe « illustration » en minuscules ou « Illustration » en majuscules en début de phrase.
L'ensemble de ces supports est répertorié dans la \emph{Liste des illustrations} du Post-Contenu.

